\begin{abstract}
Online discussions are increasingly vulnerable to coordinated narrative attacks. Attackers plant carefully framed messages to skew debate and inflame emotions, disrupting deliberation and governance. Yet existing research has largely treated such attacks as signals to detect after they appear, rather than how they spread through conversation or how interventions contain them. To move from detection to mechanism, we propose \textbf{NarrateSim},\footnote{Code and data: \url{https://github.com/SocialAI-X/NarrateSim}} a multi-agent simulation in which LLM agents reason via a Belief--Desire--Intention (BDI) process, interact over a scale-free network, and face prompt-driven attackers, enabling study of how strategic narratives take hold in discussion. Experiments reveal two patterns. On the attack side, even a few attackers move group opinion away from the no-attack baseline. Evidence- and credential-based strategies most strongly reshape stance, while victim-style storytelling raises emotional temperature. Notably, identical content can produce sharply different effects across audiences. On the defence side, \textbf{pre-bunking} outperforms reactive correction, suggesting that early inoculation is more reliable than after-the-fact repair. However, timing matters. Once emotions run high, even self-reflection prompts can amplify rather than reduce the attack. Overall, our work shifts narrative-attack research from post-hoc detection to dynamic mechanism analysis and offers practical governance strategies for defending online communities against coordinated manipulation.
\end{abstract}
